\caption{Recovery of a planted time-dependent response plane from multivariate observations. A ten-dimensional linear system contains an independently specified two-dimensional amplifier whose orientation, spectrum or excitation changes across five controlled scenarios: stationary operator, slow plane rotation, fast plane rotation, planted threshold crossing and a noise-variance jump. Six independently seeded subjects are generated per scenario; each contributes nine overlapping windows (length 80, stride 40). The M2 optimisation and M3 commutator method are estimated from $M=1$ (pale, open symbols) or $M=12$ (saturated, filled symbols) nested trajectories. (a) Median largest principal angle $\mathrm{med}\,\theta_{\mathrm{mid}}$ to the independently planted mid-window response plane. Horizontal segments connect $M=1$ and $M=12$ for the same method; short whiskers on the $M=12$ estimates denote the archived subject-cluster 95\% bootstrap intervals. (b) Median absolute difference between the inferred rate-based diagnostic and the known center-time value, $\mathrm{med}\,|\widehat R_A-R_{A,c}|$, evaluated only for admissible fits. (c) Percentage of admissible $R_A$ estimates, with 54 windows per method and scenario. The 54 windows are not independent replicates. (d) Threshold-crossing scenario: median $M=12$ estimates against the known $R_A$ at each window center. Thin vertical stems indicate the across-subject interquartile range (not confidence intervals), the horizontal dashed line denotes $R_A=1$, and the vertical dotted line marks the independently defined center-time crossing (approximately $t=117.78$). The reference $R_{A,c}$ comes from the known instantaneous planted block at the window center, not a population window-wise least-squares predictor. Window-center timestamps are retrospective and do not imply prospective threshold detection. All panels reproduce the frozen V19 T10 study without rerunning the numerical experiment.}
